% =========================================================
% Proof: Subsequences of Cauchy Sequences are Cauchy
% Source: notes/cauchy/notes-cauchy.tex
% Difficulty: Easy
% =========================================================

\subsection*{Subsequences of Cauchy Sequences are Cauchy}
\label{prf:cauchy-subsequence-cauchy}

\begin{remark*}[Return]
$\leftarrow$ \hyperref[prop:cauchy-subsequence-cauchy]{Back to Proposition in Notes}
\end{remark*}

\begin{proof}
\Claim Let $(X,d)$ be a metric space.
If $(x_n)$ is Cauchy, then every subsequence $(x_{n_k})$ is Cauchy.

\Given A Cauchy sequence $(x_n)$ and a subsequence $(x_{n_k})$.

\Goal To show $(x_{n_k})$ is Cauchy.

\Strategy Given $\varepsilon > 0$, use the Cauchy condition on $(x_n)$
to find $N$. For $j, k \ge N$, we have $n_j \ge j \ge N$ and
$n_k \ge k \ge N$, so $d(x_{n_j}, x_{n_k}) < \varepsilon$.

% -------------------------------------------------------
% YOUR PROOF HERE
% -------------------------------------------------------

\end{proof}

\begin{remark*}[Proof shape]
Observation: subsequence indices grow, so $n_k \ge k$.
The Cauchy bound on the original sequence applies directly to
any pair of subsequence terms with large enough indices.
\end{remark*}

\begin{remark*}[Dependencies]
\begin{itemize}
  \item Definition~\ref{def:cauchy}: Cauchy sequence.
  \item Fact: $n_k \ge k$ for all $k$ (subsequence indices are strictly increasing).
\end{itemize}
\end{remark*}
